% Options for packages loaded elsewhere
\PassOptionsToPackage{unicode}{hyperref}
\PassOptionsToPackage{hyphens}{url}
\documentclass[
  spanish,
]{article}
\usepackage{xcolor}
\usepackage[margin=1in]{geometry}
\usepackage{amsmath,amssymb}
\setcounter{secnumdepth}{5}
\usepackage{iftex}
\ifPDFTeX
  \usepackage[T1]{fontenc}
  \usepackage[utf8]{inputenc}
  \usepackage{textcomp} % provide euro and other symbols
\else % if luatex or xetex
  \usepackage{unicode-math} % this also loads fontspec
  \defaultfontfeatures{Scale=MatchLowercase}
  \defaultfontfeatures[\rmfamily]{Ligatures=TeX,Scale=1}
\fi
\usepackage{lmodern}
\ifPDFTeX\else
  % xetex/luatex font selection
\fi
% Use upquote if available, for straight quotes in verbatim environments
\IfFileExists{upquote.sty}{\usepackage{upquote}}{}
\IfFileExists{microtype.sty}{% use microtype if available
  \usepackage[]{microtype}
  \UseMicrotypeSet[protrusion]{basicmath} % disable protrusion for tt fonts
}{}
\makeatletter
\@ifundefined{KOMAClassName}{% if non-KOMA class
  \IfFileExists{parskip.sty}{%
    \usepackage{parskip}
  }{% else
    \setlength{\parindent}{0pt}
    \setlength{\parskip}{6pt plus 2pt minus 1pt}}
}{% if KOMA class
  \KOMAoptions{parskip=half}}
\makeatother
\usepackage{color}
\usepackage{fancyvrb}
\newcommand{\VerbBar}{|}
\newcommand{\VERB}{\Verb[commandchars=\\\{\}]}
\DefineVerbatimEnvironment{Highlighting}{Verbatim}{commandchars=\\\{\}}
% Add ',fontsize=\small' for more characters per line
\usepackage{framed}
\definecolor{shadecolor}{RGB}{248,248,248}
\newenvironment{Shaded}{\begin{snugshade}}{\end{snugshade}}
\newcommand{\AlertTok}[1]{\textcolor[rgb]{0.94,0.16,0.16}{#1}}
\newcommand{\AnnotationTok}[1]{\textcolor[rgb]{0.56,0.35,0.01}{\textbf{\textit{#1}}}}
\newcommand{\AttributeTok}[1]{\textcolor[rgb]{0.13,0.29,0.53}{#1}}
\newcommand{\BaseNTok}[1]{\textcolor[rgb]{0.00,0.00,0.81}{#1}}
\newcommand{\BuiltInTok}[1]{#1}
\newcommand{\CharTok}[1]{\textcolor[rgb]{0.31,0.60,0.02}{#1}}
\newcommand{\CommentTok}[1]{\textcolor[rgb]{0.56,0.35,0.01}{\textit{#1}}}
\newcommand{\CommentVarTok}[1]{\textcolor[rgb]{0.56,0.35,0.01}{\textbf{\textit{#1}}}}
\newcommand{\ConstantTok}[1]{\textcolor[rgb]{0.56,0.35,0.01}{#1}}
\newcommand{\ControlFlowTok}[1]{\textcolor[rgb]{0.13,0.29,0.53}{\textbf{#1}}}
\newcommand{\DataTypeTok}[1]{\textcolor[rgb]{0.13,0.29,0.53}{#1}}
\newcommand{\DecValTok}[1]{\textcolor[rgb]{0.00,0.00,0.81}{#1}}
\newcommand{\DocumentationTok}[1]{\textcolor[rgb]{0.56,0.35,0.01}{\textbf{\textit{#1}}}}
\newcommand{\ErrorTok}[1]{\textcolor[rgb]{0.64,0.00,0.00}{\textbf{#1}}}
\newcommand{\ExtensionTok}[1]{#1}
\newcommand{\FloatTok}[1]{\textcolor[rgb]{0.00,0.00,0.81}{#1}}
\newcommand{\FunctionTok}[1]{\textcolor[rgb]{0.13,0.29,0.53}{\textbf{#1}}}
\newcommand{\ImportTok}[1]{#1}
\newcommand{\InformationTok}[1]{\textcolor[rgb]{0.56,0.35,0.01}{\textbf{\textit{#1}}}}
\newcommand{\KeywordTok}[1]{\textcolor[rgb]{0.13,0.29,0.53}{\textbf{#1}}}
\newcommand{\NormalTok}[1]{#1}
\newcommand{\OperatorTok}[1]{\textcolor[rgb]{0.81,0.36,0.00}{\textbf{#1}}}
\newcommand{\OtherTok}[1]{\textcolor[rgb]{0.56,0.35,0.01}{#1}}
\newcommand{\PreprocessorTok}[1]{\textcolor[rgb]{0.56,0.35,0.01}{\textit{#1}}}
\newcommand{\RegionMarkerTok}[1]{#1}
\newcommand{\SpecialCharTok}[1]{\textcolor[rgb]{0.81,0.36,0.00}{\textbf{#1}}}
\newcommand{\SpecialStringTok}[1]{\textcolor[rgb]{0.31,0.60,0.02}{#1}}
\newcommand{\StringTok}[1]{\textcolor[rgb]{0.31,0.60,0.02}{#1}}
\newcommand{\VariableTok}[1]{\textcolor[rgb]{0.00,0.00,0.00}{#1}}
\newcommand{\VerbatimStringTok}[1]{\textcolor[rgb]{0.31,0.60,0.02}{#1}}
\newcommand{\WarningTok}[1]{\textcolor[rgb]{0.56,0.35,0.01}{\textbf{\textit{#1}}}}
\usepackage{longtable,booktabs,array}
\newcounter{none} % for unnumbered tables
\usepackage{calc} % for calculating minipage widths
% Correct order of tables after \paragraph or \subparagraph
\usepackage{etoolbox}
\makeatletter
\patchcmd\longtable{\par}{\if@noskipsec\mbox{}\fi\par}{}{}
\makeatother
% Allow footnotes in longtable head/foot
\IfFileExists{footnotehyper.sty}{\usepackage{footnotehyper}}{\usepackage{footnote}}
\makesavenoteenv{longtable}
\usepackage{graphicx}
\makeatletter
\newsavebox\pandoc@box
\newcommand*\pandocbounded[1]{% scales image to fit in text height/width
  \sbox\pandoc@box{#1}%
  \Gscale@div\@tempa{\textheight}{\dimexpr\ht\pandoc@box+\dp\pandoc@box\relax}%
  \Gscale@div\@tempb{\linewidth}{\wd\pandoc@box}%
  \ifdim\@tempb\p@<\@tempa\p@\let\@tempa\@tempb\fi% select the smaller of both
  \ifdim\@tempa\p@<\p@\scalebox{\@tempa}{\usebox\pandoc@box}%
  \else\usebox{\pandoc@box}%
  \fi%
}
% Set default figure placement to htbp
\def\fps@figure{htbp}
\makeatother
\ifLuaTeX
\usepackage[bidi=basic,shorthands=off]{babel}
\else
\usepackage[bidi=default,shorthands=off]{babel}
\fi
\ifLuaTeX
  \usepackage{selnolig} % disable illegal ligatures
\fi
\setlength{\emergencystretch}{3em} % prevent overfull lines
\providecommand{\tightlist}{%
  \setlength{\itemsep}{0pt}\setlength{\parskip}{0pt}}
\usepackage{bookmark}
\IfFileExists{xurl.sty}{\usepackage{xurl}}{} % add URL line breaks if available
\urlstyle{same}
% fallback for those not using the hyperref driver hyperxmp:
\makeatletter
\@ifundefined{xmpquote}{\newcommand{\xmpquote}[1]{#1}}{}
\makeatother
\hypersetup{
  pdftitle={Práctica 2. AED},
  pdfauthor={Javier Rodríguez Jiménez},
  pdflang={es-ES},
  hidelinks,
  pdfcreator={LaTeX via pandoc}}

\title{Práctica 2. AED}
\usepackage{etoolbox}
\makeatletter
\providecommand{\subtitle}[1]{% add subtitle to \maketitle
  \apptocmd{\@title}{\par {\large #1 \par}}{}{}
}
\makeatother
\subtitle{Introducción a la Ciencia de Datos. Máster en Ciencia de Datos- UV}
\author{Javier Rodríguez Jiménez}
\date{2026-10-26}

\begin{document}
\maketitle

{
\setcounter{tocdepth}{3}
\tableofcontents
}
Cargamos, en primer lugar, la librería

\begin{Shaded}
\begin{Highlighting}[]
\FunctionTok{library}\NormalTok{(ggplot2)}
\end{Highlighting}
\end{Shaded}

\section{Ejercicio 1}\label{ejercicio-1}

\textbf{a)} Carga el fichero weatherAUS.RData utilizando lo aprendido en la práctica 1.

\begin{Shaded}
\begin{Highlighting}[]
\NormalTok{weather\_aus }\OtherTok{\textless{}{-}} \FunctionTok{load}\NormalTok{(}\StringTok{\textquotesingle{}data/weatherAUS.RData\textquotesingle{}}\NormalTok{)}
\end{Highlighting}
\end{Shaded}

\textbf{b)} Crea un objeto ggplot denominado grafico1 en el que esté definido el conjunto de datos ds alojado en el fichero weatherAUS.RData y un aesthetics en el que en el eje \(x\) esté definida la variable MinTemp y en el eje \(y\) esté definida la variable MaxTemp.

\begin{Shaded}
\begin{Highlighting}[]
\NormalTok{grafico1 }\OtherTok{\textless{}{-}} \FunctionTok{ggplot}\NormalTok{(ds, }\FunctionTok{aes}\NormalTok{(}\AttributeTok{x =}\NormalTok{ MinTemp, }\AttributeTok{y =}\NormalTok{ MaxTemp))}
\end{Highlighting}
\end{Shaded}

\textbf{c)} Añade una capa que defina el objeto geométrico necesario para construir un gráfico de dispersión de las variables MinTemp y MaxTemp.

\begin{Shaded}
\begin{Highlighting}[]
\NormalTok{grafico1}\SpecialCharTok{+}\FunctionTok{geom\_point}\NormalTok{()}
\end{Highlighting}
\end{Shaded}

\pandocbounded{\includegraphics[keepaspectratio]{Practica2AED_JavierRodriguez_files/figure-latex/unnamed-chunk-4-1.pdf}}

\textbf{d)} Genera ahora el mismo gráfico y almacénalo en grafico2 pero define el aesthetics para visualizar las variables MinTemp y MaxTemp en los ejes de la gráfica y la variable categórica RainTomorrow como el color. A partir del objeto grafico2, representa un gráfico de dispersión en el que se muestren los puntos con el color del punto correspondiente a la variable RainTomorrow.

\begin{Shaded}
\begin{Highlighting}[]
\NormalTok{grafico2 }\OtherTok{\textless{}{-}} \FunctionTok{ggplot}\NormalTok{(ds, }\FunctionTok{aes}\NormalTok{(}\AttributeTok{x =}\NormalTok{ MinTemp, }\AttributeTok{y =}\NormalTok{ MaxTemp, }\AttributeTok{colour =}\NormalTok{ RainTomorrow))}
\NormalTok{grafico2 }\OtherTok{\textless{}{-}}\NormalTok{grafico2 }\SpecialCharTok{+} \FunctionTok{geom\_point}\NormalTok{()}
\end{Highlighting}
\end{Shaded}

\textbf{e)} Añade a grafico2 una capa para incluir un título con ggtitle() y almacénalo en grafico3.

\begin{Shaded}
\begin{Highlighting}[]
\NormalTok{grafico3 }\OtherTok{\textless{}{-}}\NormalTok{ grafico2}\SpecialCharTok{+}\FunctionTok{ggtitle}\NormalTok{(}\StringTok{\textquotesingle{}Máxima frente a mínima temperatura.\textquotesingle{}}\NormalTok{)}
\NormalTok{grafico3}
\end{Highlighting}
\end{Shaded}

\pandocbounded{\includegraphics[keepaspectratio]{Practica2AED_JavierRodriguez_files/figure-latex/unnamed-chunk-6-1.pdf}}

\textbf{f)} Almacena en formato pdf el gráfico resultante en un fichero denominado Ejercicio1.pdf

\begin{Shaded}
\begin{Highlighting}[]
\FunctionTok{ggsave}\NormalTok{(}\StringTok{\textquotesingle{}results/Ejercicio1.pdf\textquotesingle{}}\NormalTok{, grafico3)}
\end{Highlighting}
\end{Shaded}

\begin{verbatim}
## Saving 6.5 x 4.5 in image
\end{verbatim}

\section{Ejercicio 2}\label{ejercicio-2}

\textbf{a)} Utilizando ggplot haz un grafico de dispersión en el que representes las variables Temp9am y Pressure9am en los ejes \(x\) e \(y\), respectivamente.

\begin{Shaded}
\begin{Highlighting}[]
\NormalTok{grafico2\_a }\OtherTok{\textless{}{-}} \FunctionTok{ggplot}\NormalTok{(ds, }\FunctionTok{aes}\NormalTok{(}\AttributeTok{x =}\NormalTok{ Temp9am, }\AttributeTok{y =}\NormalTok{ Pressure9am))}\SpecialCharTok{+}\FunctionTok{geom\_point}\NormalTok{()}
\NormalTok{grafico2\_a}
\end{Highlighting}
\end{Shaded}

\pandocbounded{\includegraphics[keepaspectratio]{Practica2AED_JavierRodriguez_files/figure-latex/unnamed-chunk-8-1.pdf}}

\textbf{b)}Modifica el gráfico anterior para incluir información categórica proporcionada por la variable RainToday con la forma del punto.

\begin{Shaded}
\begin{Highlighting}[]
\NormalTok{grafico2\_b }\OtherTok{\textless{}{-}}\NormalTok{ grafico2\_a }\SpecialCharTok{+} \FunctionTok{aes}\NormalTok{(}\AttributeTok{shape =}\NormalTok{ RainToday)}
\NormalTok{grafico2\_b}
\end{Highlighting}
\end{Shaded}

\pandocbounded{\includegraphics[keepaspectratio]{Practica2AED_JavierRodriguez_files/figure-latex/unnamed-chunk-9-1.pdf}}

\textbf{c)} Modifica el gráfico anterior para incluir información categórica proporcionada por la variable RainToday con el color del punto. ¿Con qué representación (color o forma) se observa mejor la información que introduce la variable RainToday?

\begin{Shaded}
\begin{Highlighting}[]
\NormalTok{grafico2\_c }\OtherTok{\textless{}{-}}\NormalTok{ grafico2\_a }\SpecialCharTok{+} \FunctionTok{aes}\NormalTok{(}\AttributeTok{color =}\NormalTok{ RainToday)}
\NormalTok{grafico2\_c}
\end{Highlighting}
\end{Shaded}

\pandocbounded{\includegraphics[keepaspectratio]{Practica2AED_JavierRodriguez_files/figure-latex/unnamed-chunk-10-1.pdf}}
Se observa mejor distinguir los puntos por colores más que por forma. Al graficar todo del mismo color, las formas pueden superponerse y no distinguir los triángulos de los círculos.

\textbf{d)} Añade títulos explicativos al gráfico anterior (debe aparece un título, subtítulo, leyendas e información de las variables representadas en los ejes).

\begin{Shaded}
\begin{Highlighting}[]
\NormalTok{grafico2\_d }\OtherTok{\textless{}{-}}\NormalTok{ grafico2\_c}\SpecialCharTok{+} \FunctionTok{labs}\NormalTok{(}\AttributeTok{title =} \StringTok{"Relación entre Temperatura y Presión a las 9 am"}\NormalTok{,}
    \AttributeTok{subtitle =} \StringTok{"Distinción entre precipitación hoy."}\NormalTok{,}
    \AttributeTok{x =} \StringTok{"Temperatura @ 9am / (au)"}\NormalTok{,}
    \AttributeTok{y =} \StringTok{"Presión @ 9:00 / (au)"}\NormalTok{,}
    \AttributeTok{color =} \StringTok{"Precipitación hoy"}
\NormalTok{  )}
\end{Highlighting}
\end{Shaded}

\textbf{e)}Almacena en formato pdf el gráfico resultante en un fichero denominado Ejercicio2.pdf.

\begin{Shaded}
\begin{Highlighting}[]
\FunctionTok{ggsave}\NormalTok{(}\StringTok{\textquotesingle{}./results/Ejercicio2.pdf\textquotesingle{}}\NormalTok{, grafico2\_d)}
\end{Highlighting}
\end{Shaded}

\begin{verbatim}
## Saving 6.5 x 4.5 in image
\end{verbatim}

\section{Ejercicio 3}\label{ejercicio-3}

\textbf{a)} Carga el fichero iris\_tidy.Rdata utilizando lo aprendido en la práctica 1.

\begin{Shaded}
\begin{Highlighting}[]
\FunctionTok{load}\NormalTok{(}\StringTok{\textquotesingle{}data/iris\_tidy.Rdata\textquotesingle{}}\NormalTok{)}
\end{Highlighting}
\end{Shaded}

\textbf{b)}Utiliza la función str() para ver qué variables contiene el dataset.

\begin{Shaded}
\begin{Highlighting}[]
\FunctionTok{str}\NormalTok{(iris.tidy)}
\end{Highlighting}
\end{Shaded}

\begin{verbatim}
## 'data.frame':    600 obs. of  4 variables:
##  $ Species: Factor w/ 3 levels "setosa","versicolor",..: 1 1 1 1 1 1 1 1 1 1 ...
##  $ Part   : Factor w/ 2 levels "Petal","Sepal": 2 2 2 2 2 2 2 2 2 2 ...
##  $ Measure: chr  "Length" "Length" "Length" "Length" ...
##  $ Value  : num  5.1 4.9 4.7 4.6 5 5.4 4.6 5 4.4 4.9 ...
\end{verbatim}

\textbf{c)}Crea un objeto ggplot denominado iris.grafico1 en el que esté definido el conjunto de datos iris.tidy alojado en el fichero iris.tidy.RData y un aesthetics en el que en el eje \(x\) esté definida la variable Species y en el eje \(y\) esté definida la variable Value.

\begin{Shaded}
\begin{Highlighting}[]
\NormalTok{iris.grafico1 }\OtherTok{\textless{}{-}} \FunctionTok{ggplot}\NormalTok{(iris.tidy, }\FunctionTok{aes}\NormalTok{(}\AttributeTok{x =}\NormalTok{ Species, }\AttributeTok{y =}\NormalTok{ Value))}
\end{Highlighting}
\end{Shaded}

\textbf{d)} Añade una capa que defina el objeto geométrico necesario para construir un gráfico de dispersión de las variables Species y Value.

\begin{Shaded}
\begin{Highlighting}[]
\NormalTok{iris.grafico1}\SpecialCharTok{+}\FunctionTok{geom\_point}\NormalTok{()}
\end{Highlighting}
\end{Shaded}

\pandocbounded{\includegraphics[keepaspectratio]{Practica2AED_JavierRodriguez_files/figure-latex/unnamed-chunk-16-1.pdf}}

\textbf{e)} Modifica el aesthetics de iris.grafico1 para visualizar la variable categórica Part y almacénalo en iris.grafico2. A partir del objeto iris.grafico2, representa un gráfico de dispersión en el que se muestren los puntos dados por las variables Species y Value, y el color del punto se corresponda con el valor de la variable Part.

\begin{Shaded}
\begin{Highlighting}[]
\NormalTok{iris.grafico2 }\OtherTok{\textless{}{-}}\NormalTok{ iris.grafico1}\SpecialCharTok{+}\FunctionTok{geom\_point}\NormalTok{()}\SpecialCharTok{+}\FunctionTok{aes}\NormalTok{(}\AttributeTok{color=}\NormalTok{Part)}
\NormalTok{iris.grafico2}
\end{Highlighting}
\end{Shaded}

\pandocbounded{\includegraphics[keepaspectratio]{Practica2AED_JavierRodriguez_files/figure-latex/unnamed-chunk-17-1.pdf}}

\textbf{f)} Incluye información de la variable Measure haciendo dos gráficos de dispersión (dispuestos en filas) con facet\_grid(), como en el apartado de introducción, para los dos valores que toma la variable Measure. Prueba a cambiar la disposición de los gráficos para que sea en columnas ¿Qué disposición, en filas o en columnas, permite comparar mejor los valores de las dos variables?

\begin{Shaded}
\begin{Highlighting}[]
\NormalTok{iris.grafico2 }\SpecialCharTok{+} \FunctionTok{facet\_grid}\NormalTok{(}\AttributeTok{rows =} \FunctionTok{vars}\NormalTok{(Measure))}
\end{Highlighting}
\end{Shaded}

\pandocbounded{\includegraphics[keepaspectratio]{Practica2AED_JavierRodriguez_files/figure-latex/unnamed-chunk-18-1.pdf}}

\begin{Shaded}
\begin{Highlighting}[]
\NormalTok{iris.grafico2 }\SpecialCharTok{+} \FunctionTok{facet\_grid}\NormalTok{(}\AttributeTok{cols =} \FunctionTok{vars}\NormalTok{(Measure))}
\end{Highlighting}
\end{Shaded}

\pandocbounded{\includegraphics[keepaspectratio]{Practica2AED_JavierRodriguez_files/figure-latex/unnamed-chunk-18-2.pdf}}
La mejor disposición es aquella que alinea la variable Value, es decir, en la que obtenemos 1 fila y 2 gráficas. Pese a que la configuración análoga puede ser útil puesto que nos alinea en una columna el tipo de especies, no podemos comparar directamente la variable del eje y que es realmente nuestro objetivo.

\textbf{g)}Cambia ahora la capa geom\_point() del gráfico anterior por la capa geom\_jitter(). ¿Qué diferencia observas?

\begin{Shaded}
\begin{Highlighting}[]
\NormalTok{iris.grafico2 }\SpecialCharTok{+} \FunctionTok{geom\_jitter}\NormalTok{()}\SpecialCharTok{+} \FunctionTok{facet\_grid}\NormalTok{(}\AttributeTok{rows =} \FunctionTok{vars}\NormalTok{(Measure))}
\end{Highlighting}
\end{Shaded}

\pandocbounded{\includegraphics[keepaspectratio]{Practica2AED_JavierRodriguez_files/figure-latex/unnamed-chunk-19-1.pdf}}

\begin{Shaded}
\begin{Highlighting}[]
\NormalTok{iris.grafico2 }\SpecialCharTok{+} \FunctionTok{geom\_jitter}\NormalTok{()}\SpecialCharTok{+} \FunctionTok{facet\_grid}\NormalTok{(}\AttributeTok{cols =} \FunctionTok{vars}\NormalTok{(Measure))}
\end{Highlighting}
\end{Shaded}

\pandocbounded{\includegraphics[keepaspectratio]{Practica2AED_JavierRodriguez_files/figure-latex/unnamed-chunk-19-2.pdf}}
La diferencia principal es que ahora los datos en vez de estar perfectamente alineados sobre cada una de la categoría de especies se ha añadido un ruido o una pequeña variación en la posición de nuestros datos. Esto es de gran utilidad en nuestro caso porque encontrábamos todos los datos apilados sobre el eje \(x\) y, realmente, no se podía apreciar la cantidad de datos que repetían valor.

\textbf{h)}Almacena en formato pdf el gráfico resultante en un fichero denominado Ejercicio3.pdf.

\begin{Shaded}
\begin{Highlighting}[]
\FunctionTok{ggsave}\NormalTok{(}\StringTok{\textquotesingle{}./results/Ejercicio3.pdf\textquotesingle{}}\NormalTok{)}
\end{Highlighting}
\end{Shaded}

\begin{verbatim}
## Saving 6.5 x 4.5 in image
\end{verbatim}

\section{Ejercicio 4}\label{ejercicio-4}

Carga el fichero Ejercicio4\_Data.Rdata. Examina las variables que tiene dicho dataset.
Representa un gráfico de puntos, en el que se muestre la fecha en el eje 𝑥
y las ventas en el eje 𝑦
. Queremos que los datos de cada fabricante sean de un color distinto. ¿Es este gráfico válido para evaluar la evolución de las ventas de cada fabricante?
Representa un gráfico como el anterior en el que se muestre la línea que une los puntos que antes has representado.
Modifica el gráfico anterior para que la tendencia de cada fabricante este trazada con un tipo de línea diferente.
Añade el siguiente título al gráfico ``Ventas mensuales de diferentes compañías tecnológicas''.
Almacena en formato pdf el gráfico resultante en un fichero denominado Ejercicio4.pdf.

\end{document}
